%! AP Statistics — Unit 8, Lesson 8.2 — Homework
\documentclass[10pt]{article}
\usepackage{apstats-article}
\usepackage{apstats-boxes}

\begin{document}

\pageheader{Unit 8, Lesson 8.2}{Homework}
\namedateperiod

\vspace{0.1in}
\textbf{Directions:} For each problem, set up a chi-square goodness of fit test as directed.
Show all work clearly.

\begin{enumerate}

  \item \textbf{Genetics.} A genetics researcher expects offspring to appear in a 9:3:3:1
        phenotype ratio. In a study, she observes 160 offspring with counts: Phenotype A = 92,
        Phenotype B = 28, Phenotype C = 26, Phenotype D = 14.

        \begin{enumerate}[label=\alph*.]
          \item State H$_0$ and H$_a$ for a chi-square goodness of fit test.

                H$_0$: \writelines{2}

                H$_a$: \writeline

          \item Calculate the expected count for each phenotype. (\textit{Hint:}
                Total ratio parts $= 9 + 3 + 3 + 1 = 16$.)

                \begin{center}
                \renewcommand{\arraystretch}{1.4}
                \begin{tabular}{lcccc}
                    \textbf{Phenotype} & A & B & C & D \\
                    \hline
                    Null proportion & $9/16$ & $3/16$ & $3/16$ & $1/16$ \\
                    \hline
                    Expected count & \blank{2cm} & \blank{2cm} & \blank{2cm} & \blank{2cm} \\
                \end{tabular}
                \end{center}

          \item Is the Large Counts condition met? Justify.
                \writelines{1}
        \end{enumerate}

  \item \textbf{Transportation Survey.} A local government claims that commuters use four
        modes of transportation in equal proportions (25\% each). In a random sample of
        200 commuters: car = 72, bus = 58, bike = 36, walk = 34.

        \begin{enumerate}[label=\alph*.]
          \item State H$_0$ and H$_a$ for a test of whether the proportions are equal.

                H$_0$: \writeline

                H$_a$: \writeline

          \item Calculate the expected count for each transportation mode.

                $E = n \cdot p_0 = $ \blank{6cm}

                \begin{center}
                \renewcommand{\arraystretch}{1.3}
                \begin{tabular}{lcccc}
                    Mode & Car & Bus & Bike & Walk \\
                    \hline
                    Expected & \blank{1.8cm} & \blank{1.8cm} & \blank{1.8cm} & \blank{1.8cm} \\
                \end{tabular}
                \end{center}

          \item Check both conditions for a chi-square GOF test.
                \begin{itemize}
                    \item \textbf{Random:} \blank{7cm}
                    \item \textbf{Large Counts:} \blank{7cm}
                \end{itemize}
        \end{enumerate}

  \item \textbf{True or False.} Circle your answer and explain.

        ``If the sample size is large enough, we only need to check that each
        \emph{observed} count is greater than 5 for the chi-square GOF test.''

        \begin{center}
        \textbf{True} \qquad \textbf{False}
        \end{center}

        Explanation: \writelines{3}

\end{enumerate}

\end{document}
